%!TEX root = ./../../../springer-mono-format/book.tex
%!TEX root = ./../../../thesis-trier.tex

In general, the computational complexity might increase or decrease
when concatenating languages.
For example, 
by \cref{constr_sync:thm:introductory_examples},
\cref{gen_results:cor:factors_etc_complexity}
and as $\fact{(a+b+c)^*c(a+b)^*} = \fact{(a+b)^*}$,
$\constrsync{(a+b)^*}$, $\constrsync{c(a+b)^*}$,
and $\constrsync{(a+b+c)^*c(a+b)^*}$ are 
complete for \NL, \PSPACE, and \NL, respectively.
Here, we identify conditions on concatenating languages
to obtain complexity lower bounds.
First, we give two results
for concatenating a language
to the front of a constraint language such that hardness is preserved, where
\cref{constrsync:prop:UanL} is related to languages contained in a subalphabet
and \cref{constr_sync:thm:prefix_before} to prefix codes.
Then, \cref{constr_sync:prop:Gamma_suffix}
is concerned with concatenation to the back.
The conditions presented here are used to deduce hardness results in later sections.

The reductions of the present section involve operations
as mentioned in \cref{prelim:rem:reductionsAC0}
and can be realized by $\AC^0$ circuits.

%\subsubsection{Concatenating to the Front}

\begin{proposition}
	\label{constrsync:prop:UanL}
	Let $L \subseteq \Sigma^*$ 
	and $U \subseteq \Gamma^*$ for some $\Gamma \subseteq \Sigma$
	be regular languages and $w \in \Sigma^*$.
	Suppose one of the following holds:
	\begin{enumerate}
	\item $L \cap \Gamma\Sigma^* = \emptyset$ with $w \in \Gamma^*$, or
	
	\item $w \notin \Gamma^*$.
	\end{enumerate}
	Then $\constrsync{wL} \leACZEROm \constrsync{UwL}$.
\end{proposition}         
\begin{proof} 
	First, assume $L \cap \Gamma\Sigma^* = \emptyset$ and $w \in \Gamma^*$.	
	Let $\semiautomat{A}{Q}{\Sigma}{\delta}$. We construct a semiautomaton $\mathcal A'$ that has a synchronizing
	word in $UwL$ if and only if $\mathcal A$ has a synchronizing word in $wL$.
	
	\medskip
	\noindent\underline{Construction:} Let $\semiautomat{A'}{Q \times \{0,1\}}{\Sigma}{\delta'}$
	where, for $i \in \{0,1\}$, $q \in Q$ and $x \in \Sigma$,
	\[
	\applytrans{\delta'}{(q, i)}{x} = 
	\begin{cases}
		(q, 0) & \text{if } i = 0 \text{ and } x \in \Gamma, \\
		(\applytrans{\delta}{q}{wx}, 1) & \text{if } i = 0 \text{ and } x \notin \Gamma, \\
		(\applytrans{\delta}{q}{x}, 1) & \text{if } i = 1.
	\end{cases}
	\]
	
	\noindent\underline{Properties:} Let $q \in Q$, $u \in U$ and $v \in L$.
	Then, as $v \notin \Gamma\Sigma^*$ and $uw \in \Gamma^*$,
	we have 
	\begin{align}
		\applytrans{\delta'}{(q, 0)}{uwv} & = \applytrans{\delta'}{(q, 0)}{v}, \label{constrsync:eqn:UL1} \\
		\applytrans{\delta'}{(q, 0)}{v}     & = (\applytrans{\delta}{q}{wv}, 1), \label{constrsync:eqn:UL3} \\
		\applytrans{\delta'}{(q, 1)}{uwv} & = (\applytrans{\delta}{q}{uwv}, 1). \label{constrsync:eqn:UL2}
	\end{align} 
	
	\noindent\underline{Verification:} Let $u \in U$ and $v \in L$.
	Then we have
	\begin{align*}
		\applytrans{\delta'}{Q \times \{0,1\}}{uwv} 
		& = \applytrans{\delta'}{Q \times \{0\}}{uwv} \cup \applytrans{\delta'}{Q \times \{1\}}{uwv} \\
		& = \applytrans{\delta'}{Q \times \{0\}}{v} \cup \applytrans{\delta'}{Q \times \{1\}}{uwv} & \tag{Eqn.~\eqref{constrsync:eqn:UL1}} \\
		& = \delta(Q, wv) \times \{1\} \cup \delta'(Q \times \{1\}, uwv) & \tag{Eqn.~\eqref{constrsync:eqn:UL3}} \\
		& = \delta(Q, wv) \times \{1\} \cup \delta(Q, uwv) \times \{1\} & \tag{Eqn.~\eqref{constrsync:eqn:UL2}} \\
		& = \delta(Q, wv)  \times \{1\}.                                   & \tag{$\delta(\delta(Q, u), wv) \subseteq  \delta(Q, wv)$}
	\end{align*}
	So $|\applytrans{\delta'}{Q \times \{0,1\}}{uwv}| = 1$
	is equivalent to $|\delta(Q,wv)| = 1$
	and $\mathcal A'$ has a synchronizing word in $UwL$
	if and only if $\mathcal A$ has a synchronizing word in $wL$.
	
    \medskip 
    
    Lastly, for the remaining case, assume $w \notin \Gamma^*$.
    Then $w = ubv$  for some $u \in \Gamma^*$, $b \notin \Gamma$, and $v \in \Sigma^*$.
    Set $L' = bvL$. We have $L' \cap \Gamma\Sigma^* = \emptyset$. 
    So, applying the previous case to $L', u$ and $U$,
    we get 
    $\constrsync{uL'} \leACZEROm \constrsync{UuL'}$
    where $uL' = wL$.
\end{proof}             

\begin{example}\label{constrsync:ex:PSPACEhard}
	By~\crefitem{constr_sync:thm:introductory_examples}{constrsync:ex:item:abcL},
	$\constrsync{c(a+b)^*}$	is \PSPACE-complete.
	By \cref{constrsync:prop:UanL}, the problems
	$\constrsync{aabbc(a+b)^*}$ and $\constrsync{a^*c(a+b)^*}$
	are \PSPACE-complete as well.
\end{example}

For reference, we list two special cases  of~\cref{constrsync:prop:UanL}.

\begin{corollary}
	\label{constrsync:cor:UanL}
	Let $L \subseteq \Sigma^*$ be regular.
	
	\begin{enumerate}
		\item \label{constrsync:cor:UanL:1} If $L \cap \Gamma\Sigma^* = \emptyset$,
		$\Gamma \subseteq \Sigma$ and $U \subseteq \Gamma^*$ is regular, then $\constrsync{L} \leACZEROm \constrsync{UL}$. 
		
		\item \label{constrsync:cor:UanL:2} $\constrsync{wL} \leACZEROm \constrsync{a^*wL}$	for $w \notin a^*$ and $a \in \Sigma$.
	\end{enumerate}
\end{corollary}                    

Next, we present a condition related to prefix codes.

\begin{theorem}
	\label{constr_sync:thm:prefix_before}
	\label{constr_sync:thm:in_font}
	%Let $U, V \subseteq \Sigma^*$ be regular and $P = U \setminus U\Sigma^+$ (those words in $U$ 
	%without a proper prefix in $U$).
	%Then $\constrsync{V} \le_m^{\log} \constrsync{PV}$. 	
	Let $L \subseteq \Sigma^*$ be regular 
	and $P \subseteq \Sigma^*$ be a prefix code that is a regular language. 
	Then 
	$
	\constrsync{L} \leACZEROm \constrsync{PL}.
	$	
\end{theorem}
\begin{proof}    	
	By~\crefitem{constr_sync:prop:aut_prefix_free}{constr_sync:prop:aut_prefix_free:item:trim} 
	there exists a trim PDFA
	$\automat{C}{R}{\Sigma}{\rho}{r_0}{\{e\}}$
	such that $P = L(\mathcal C)$
	and where the state~$e$ has no outgoing transitions.
	Suppose $\semiautomat{A}{Q}{\Sigma}{\delta}$ is an input semiautomaton. 
	
	
	% bild oder beispiel?
	\medskip 	
	\noindent\underline{Construction:} Let $\semiautomat{A'}{R \times Q}{\Sigma}{\delta'}$
	where, for $(r,q) \in R \times Q$ and $x \in \Sigma$,
	\[
	\applytrans{\delta'}{(r, q)}{x} = \begin{cases}
		(\applytrans{\rho}{r}{x}, q)   & \text{if } \applytrans{\rho}{p}{x} \text{ is defined,} \\
		(e, \applytrans{\delta}{q}{x}) & \text{if } \rho(p, x) \text{ is undefined. }
	\end{cases} 
	\]
	
	
	
	 %Property~\ref{constrsync:prefix:prop:e} readily follows from the definition of $\delta'$.
	%For the other properties, we give the proofs.  
	\medskip 	
	\noindent\underline{Properties:} Property~\ref{constrsync:prefix:defdeltastar} follows readily by induction.
	We give the proof for Property~\ref{constrsync:prefix:prop:exactsync}.
	\begin{enumerate}    	
		
		\item \label{constrsync:prefix:defdeltastar} For $r \in R$, $q \in Q$ and $u \in \Sigma^*$ we have
		\begin{equation}\label{constrsync:prefix:eqn:transitionword}
			\delta'((r, q), u)
		=
		\begin{cases}
		(e, \delta(q, xu_2)) & \text{if } \exists u_1, u_2 \in \Sigma^* \ \exists x \in \Sigma \colon u = u_1 x u_2 \land \\ 
		                    & \qquad \rho(r, u_1) \text{ is defined } \land \\ 
		                    & \qquad \rho(r, u_1 x) \text{ is undefined}, \\
		(\rho(r, u), q)     & \text{if } \rho(r, u) \text{ is defined. } 
		\end{cases}
		\end{equation}
				
\begin{comment} 				
		\item \label{constrsync:prefix:defdeltastar}
		For $r \in R$, $q \in Q$ and $u \in \Sigma^*$ we have	
		\[
		\delta'((r, q), u)
		=
		\begin{cases}
		(e, \delta(q, u)) & \text{if } r = e, \\
		(\rho(r, u), q) & \text{if } r \ne e \text{ and } \rho(r, u) \text{ is defined. }
		\end{cases}
		\]
		
		\begin{proofitem}{\ref{constrsync:prefix:defdeltastar}}
		For $r = e$ the claim follows readily from the definition of $\delta'$
		and the fact that $\rho(r, x)$ is undefined for each $x \in \Sigma$.
		Next, suppose $r \ne e$ and $\rho(r,u)$ is defined. 
		We do induction on $u$. For $u = \varepsilon$
		the claim is clear. So, suppose we have $u = vx$
		with $v \in \Sigma^*$ and $x \in \Sigma$.
		Then $\rho(r, v)$ is defined and, by the definition of $\delta'$,
		we obtain $\delta'((\rho(r, v), q), x) = (\rho(r, vx), q)$
		and the claim follows by induction.		
		\end{proofitem} 
\end{comment}
\begin{comment} 
		\item \label{constrsync:prefix:prop:e} Let $u \in \Sigma^*$ and $q \in Q$. Then $\applytrans{\delta'}{(e, q)}{u} = (e, \applytrans{\delta}{q}{u})$.
		
		\item  \label{constrsync:prefix:prop:prefixP}  Let $u \in \pref{P} \setminus P$ and $q \in Q$.
		Then $\rho(r_0, u)$ is defined and $\applytrans{\delta'}{(r_0, q)}{u} = (\rho(r_0, u), q)$
		with $\rho(r_0, u) \ne e$.
		
		\begin{proofitem}{\ref{constrsync:prefix:prop:prefixP}}
		We do induction on $u$. For $u = \varepsilon$ 
		the claim is clear. So, suppose we have $u = u'x$ for some $u' \in \Sigma^*$ and $x \in \Sigma$.
		Then, $u' \in \pref{P}\setminus P$ and, by induction hypothesis, 
		$\rho(r_0, u')$
		is defined, $\delta'((r_0, q), u') = (\rho(r_0, u'), q)$ and $\rho(r_0, u') \ne e$.    	 
		So, as $P$ is a prefix code and $u$ a proper prefix,
		$\rho(r_0, u'x)$ is defined with $\rho(r_0, u'x) \ne e$
		and, by the definition of $\delta'$,
		we find 
		\[ 
		\delta'((r_0, q), u) = \delta'((\rho(r_0, u'), q), x) = (\rho(r_0, u'x), q),
		\]
		which gives the claim by induction on the structure of $u$.
		\end{proofitem}
		
		\item \label{constrsync:prefix:prop:wordP} Let $u \in P$ and $q \in Q$. Then $\applytrans{\delta'}{(r_0, q)}{u} = (e, q)$.
		
		\begin{proofitem}{\ref{constrsync:prefix:prop:wordP}}
		Let $u = u'x$ for $u' \in \pref{P}$ and $x \in \Sigma$.
		By Property~\ref{constrsync:prefix:prop:prefixP}
		we get $\applytrans{\delta'}{(r_0, q)}{u'} = (\rho(r_0,u'), q)$
		with $\rho(r_0, u') \ne e$.
		Then, by the definition of $\delta'$, this yields
		$\applytrans{\delta'}{(r_0, q)}{u} = \applytrans{\delta'}{(\applytrans{\rho}{r_0}{u'}, q)}{x}
		= (\applytrans{\rho}{\applytrans{\rho}{r_0}{u'}}{x}, q) = (e, q)$.
		\end{proofitem}
\end{comment}
\begin{comment} 
		\item \label{constrsync:prefix:prop:undefined} Let $w \in \Sigma^*$ and $r \in R$. If $\rho(r, w)$ is undefined,
		then $\delta'((r, q), w) \in \{e\} \times Q$ for each $q \in Q$.
		
		\begin{proofitem}{\ref{constrsync:prefix:prop:undefined}}
		Let $w = w_1 x w_2$ with $w_1, w_2 \in \Sigma^*$, $x \in \Sigma$ and 
		$w_1$ of maximal length
		such that $\rho(r, w_1)$ is defined.
		By Property~\ref{constrsync:prefix:defdeltastar} 
		we have $\delta'((r, q), w_1) = (\rho(r, w_1), q)$.
		Then, by the definition of $\delta'$, as $\rho(\rho(r, w_1), x)$ is undefined,
		we find 
		\[
		\delta'((\rho(r, w_1), q), x) = (e, \delta(q, x)).
		\]
		By Property~\ref{constrsync:prefix:defdeltastar} 
		we obtain $\delta'((e, \delta(q, x)), w_2) = (e, \delta(q, xw_2)) \in \{e\} \times Q$.
		\end{proofitem}
\end{comment}
		
		\item \label{constrsync:prefix:prop:exactsync} Let $w \in \Sigma^*$ be such that $\applytrans{\rho}{R}{w} = \{e\}$.
		Then $\applytrans{\delta'}{R \times Q}{w} = \{e\} \times Q$.
		
		
		\Cref{constrsync:prefix:eqn:transitionword} 
		gives $\delta'(R \times Q, w) \subseteq \{e\} \times Q$.
		Furthermore, there exists $r \in R$ such that $\rho(r, w) = e$.
		Then, for each $q \in Q$, we have $\delta'((r, q), w) = (e,q)$
		and this gives $\{e\} \times Q = \delta'(\{r\} \times Q, w) \subseteq \delta'(R \times Q, w)$. 	    	
	\end{enumerate}   
	
	\medskip 
	\noindent\underline{Verification:} 
	By Property~\ref{constrsync:prefix:defdeltastar}
	we have $\delta'(\{r_0\} \times Q, uv) = \{e\} \times \delta(Q, v)$
	for $u \in P$ and $v \in \Sigma^*$. (note that, since~$P$ is a 
	prefix code, $uv$ can only be factored into a word from~$P$
	and an arbitrary word in one way)
	So, $|\delta'(R \times Q, uv)| = 1$
	for $u \in P$ and $v \in L$
	implies $|\delta(Q, v)| = 1$.	
	
	Conversely, suppose $|\delta(Q, v)| = 1$
	for $v \in L$. 
	By Property~\ref{constrsync:prefix:defdeltastar} and as $\rho(e,x)$ is undefined for each $x \in \Sigma$
	we have $\delta'(\{e\} \times Q, v) = \{e\} \times \delta(Q, v)$ (here, we used that~$P$ is a prefix code).
	By \crefitem{constr_sync:prop:aut_prefix_free_exact_sync}{constr_sync:prop:aut_prefix_free:item:sync} 
	there exists $w \in P$ with
	$\rho(R, w) = \{e\}$.
	Then, by Property~\ref{constrsync:prefix:prop:exactsync} 
	we have $\applytrans{\delta'}{R \times Q}{w} \subseteq \{e\} \times Q$ (note that
	we might have that $\rho(r, w)$ is undefined for some $r \in R$ 
	and that, in fact, we have equality since $\delta'(\{r_0\}\times Q), w) = \{e\} \times Q$).
	So $|\delta'(R \times Q, wv)| = 1$.
\end{proof} 

\begin{remark}
	For $L = \{c\}$ and $P = \{a,b\}^*c$
	we have that $\constrsync{L}$ 
	and $\constrsync{PL}$
	are in $\AC^0$
	and \PSPACE-complete, respectively,
	by \cref{constr_sync:thm:introductory_examples}.
	However, the complexity does not have to increase when concatenating
	a prefix code that is regular to the front of a constraint
	language as can be seen with $L = \Sigma^*$.
	Then $\constrsync{L} \eqACZEROm \constrsync{PL}$
	by \cref{constrsync:genresults:cor:left-right-ideal}.		
\end{remark}

As $\{u\}$ for $u \in \Sigma^*$ is a prefix code, we obtain the following corollary.

\begin{corollary}
	Let $L \subseteq \Sigma^*$ be regular and $u \in \Sigma^*$.
	Then $\constrsync{L} \leACZEROm \constrsync{uL}$.
\end{corollary}

%\subsubsection{Concatenating to the Back} 

Next, we state a result on concatenating
languages to the back of the constraint language.

\begin{proposition}
	\label{constr_sync:prop:Gamma_suffix}
	Let $L \subseteq \Sigma^*$ be regular and $\Gamma \subseteq \Sigma$
	such that $L \cap \Sigma^* \Gamma = \emptyset$.
	If $U \subseteq \Gamma^*$ is a regular non-empty language,
	then $\constrsync{L} \leACZEROm \constrsync{LU}$.   
\end{proposition}
\begin{proof}
	We can assume $L \ne \emptyset$. 
	Let $\semiautomat{A}{Q}{\Sigma}{\delta}$. We construct a semiautomaton $\mathcal A'$ that has a synchronizing
	word in $LU$ if and only if $\mathcal A$ has a synchronizing word in $L$.
	\FIG~\ref{fig:irred-reduction} shows the basic idea of the construction graphically.
	\begin{comment} 
		For each $q \in Q$ and $a \in \Gamma$, set $S_{q,a} = \{ \applytrans{\delta}{q}{a^i} \mid i \ge 0 \}$.
		Define  
		$\semiautomat{A'}{Q'}{\Sigma}{\delta'}$ with
		\[
		Q' = \bigcup_{q \in Q} \bigcup_{a \in \Gamma} \{q\} \times S_{q,a}.
		\]
		and, for $(q, p) \in \{q\} \times S_{q,a}$, $x \in \Sigma$
		and $p = \applytrans{\delta}{q}{a^i}$, \todo{repräsentantenunabhgängig zeigen.}
		\[
		\applytrans{\delta'}{(q, p)}{x} = \begin{cases}
			(q, \applytrans{\delta}{p}{a})     & \text{if } x = a \\
			(\applytrans{\delta}{q}{a^ix}, \applytrans{\delta}{q}{a^ix})  & \text{if } x \ne a.
		\end{cases}
		\]
		
		Let $uv$ with $u \in L$ and $v \in U$ be a synchronizing word for $\mathcal A'$.
		Note that $u$ does not end with a letter in $\Gamma$.
		So, if $|u| > 0$,\todo{genauer?}
		we have $\applytrans{\delta'}{Q'}{u}\subseteq \{ (q,q) \mid q \in Q \}$.
		For $(q,q)$ we have $\applytrans{\delta'}{(q,q)}{a^i} = (q, \applytrans{\delta}{q}{a^i})$
		for all $i \ge 0$.
		Hence, $|\applytrans{\delta'}{Q'}{u}| = |\applytrans{\delta'}{(q,q)}{v}|$
	\end{comment} 
	
	\begin{figure}[t]
		\centering
		{			
			\begin{tikzpicture}[rounded corners=10pt]						
				% Main automaton
				\node[yshift=1cm] (A) at (0,0)
				{$\mathcal A'_{|\{(q,q)\mid q \in Q\},\,\Sigma\setminus\Gamma}
					\cong
					\mathcal A_{|\Sigma\setminus\Gamma}$};
				
				\node[draw, dashed, rectangle, fit={(-5,1.4) (5,-1.4)}] {};
				
				% Left column
				\node (q) at (-3.5,-1) {$(p,p)$};				
				\node[draw,circle,inner sep=2pt] (dummyq) at (-3.5,-2.5) {};
				\node at (-3.5,-2.75) {$\vdots$};				
				\node (labSq) at (-3.5,-3.7) {$\{p\}\times Q$};
				\node[draw,dashed,rectangle,inner sep=14pt, fit={(q) (labSq)}] (Sq) {};
				
				% Right column
				\node (p) at (3.5,-1) {$(q,q)$};				
				\node[draw,circle,inner sep=2pt] (dummyp) at (3.5,-2.5) {};
				\node at (3.5,-2.75) {$\vdots$};				
				\node (labSp) at (3.5,-3.7)	{$\{q\}\times Q$};
				\node[draw,dashed,rectangle,inner sep=14pt,	fit={(p) (labSp)}] (Sp) {};
				
				% Target states
				\node[draw,circle,inner sep=2pt] (r)  at (-1.3,-0.4) {};				
				\node[draw,circle,inner sep=2pt] (r2) at (1.3, -0.4) {};
				\node[above=0.15cm of r,  xshift=-0.4cm] {$(\delta(p,ua),\delta(p,ua))$};
				\node[above=0.15cm of r2, xshift=+0.4cm] {$(\delta(q,ua),\delta(q,ua))$};
				
				% Arrows
				\path[->] (q) edge[left,pos=.7] node {$u\in\Gamma^*$} (dummyq);				
				\path[->] (p) edge[right,pos=.7] node {$u\in\Gamma^*$} (dummyp);
				
				\path[->] (dummyq) edge[bend right,below,pos=.4] node [fill=white, anchor=center] 
				{$a \in \Sigma \setminus \Gamma$} (r);
				\path[->] (dummyp) edge[bend left,below,pos=.4]  node [fill=white, anchor=center] 
				{$a \in \Sigma \setminus \Gamma$} (r2);			
			\end{tikzpicture}
					
		}
		\caption[Schematic of reduction from the proof of Proposition~\ref{constr_sync:prop:Gamma_suffix}]{Schematic of the reduction from the proof of Proposition~\ref{constr_sync:prop:Gamma_suffix}.
			For a state $q$ and word $u \in \Gamma^*$ we ``branch out'' into a copy $\{q\} \times Q$.
			For a letter from $a \in \Sigma \setminus \Gamma$,
			we return to the state $(\applytrans{\delta}{q}{ua}, \applytrans{\delta}{q}{ua})$.
%			corresponding to the state in $\mathcal A$ when reading~$ua$.
			This is indicated 
			for two states $p$ and $q$.
		}
		\label{fig:irred-reduction}
	\end{figure}
	
	
	\medskip
	\noindent\underline{Construction:} Let 
	$\semiautomat{A'}{Q'}{\Sigma}{\delta'}$
	where\footnote{Actually, setting
		\[
		Q' = \bigcup_{q \in Q} \{q\} \times S_q
		\]
		where $S_{q} = \{ \applytrans{\delta}{q}{u} \mid u \in \Gamma^* \}$
		for $q \in Q$
		would be sufficient. However, it is not clear whether the
		sets $S_q$ can be computed by an \LOGSPACE-uniform $\AC^0$ circuit family
		whereas $Q \times Q$
		can be computed by such a circuit family.} 
	$Q' = Q\times Q$
	and, for $(q, p) \in \{q\} \times Q$, $x \in \Sigma$,
	\[
	\applytrans{\delta'}{(q, p)}{x} = \begin{cases}
		(q, \applytrans{\delta}{p}{x})     & \text{if } x \in \Gamma, \\
		(\applytrans{\delta}{p}{x}, \applytrans{\delta}{p}{x})  & \text{if } x \notin \Gamma.
	\end{cases}
	\]
		
	\noindent\underline{Properties:} Next, we collect some properties of the given construction.	
	\begin{enumerate} 
		
		\item \label{constr_sync:item:word_in_Gamma} For $w \in \Gamma^*$ and $q,p \in Q$, we have
		$
		%\begin{equation}\label{constr_sync:eqn:word_in_Gamma}
		\applytrans{\delta'}{(q,p)}{w}
		= (q, \applytrans{\delta}{p}{w}).
		%\end{equation} 
		$
		
		\begin{proofitem}{\ref{constr_sync:item:word_in_Gamma}} 
		This is clear for $w = \varepsilon$.
		Next, let $w = ux$ for some $x \in \Sigma$ and $u \in \Sigma^*$.
		Inductively, we have $\applytrans{\delta'}{(q,p)}{ux}
		= \applytrans{\delta'}{(q,\delta(p, u))}{x}$.
		Since $x \in \Gamma$,
		by the definition of the transition function
		we find $\applytrans{\delta'}{(q,\delta(p, u))}{x} = (q, \delta(p,w))$.			
		\end{proofitem}		
		
		\item \label{constr_sync:item:Gamma_suffix} If $w \in \Sigma^* (\Sigma \setminus \Gamma)$ and $q,p \in Q$, then
		$
		%\begin{equation}\label{constr_sync:eqn:Gamma_suffix}
		\applytrans{\delta'}{(q,p)}{w} = (\applytrans{\delta}{p}{w}, \applytrans{\delta}{p}{w}).
		%
		%\end{equation}
		$
				
		\begin{proofitem}{\ref{constr_sync:item:Gamma_suffix}} 
		By assumption, we can write $w = w_0 u_1 w_1 \cdots u_n w_n$ for some $n \ge 0$
		with $w_0, w_1, \ldots, w_n \in (\Sigma \setminus \Gamma)^+$
		and $u_1, \ldots, u_n \in \Gamma^*$. For $n = 0$,
		the claim follows by an easy induction from the defining equation of~$\delta'$.
		Next, suppose $n \ge 1$. Since $|w_{n-1}| \ge 1$,
		inductively, we obtain
		\[
		\delta'((q,p), w') 
		= (\delta(p, w'), \delta(p, w'))
		\]
		for $w' = w_0 u_1 w_1 \cdots w_{n-1}$. By Property~\ref{constr_sync:item:word_in_Gamma}
		we have
		\[
		\delta'((\delta(p, w'), \delta(p, w')), u_n) = (\delta(p, w'), \delta(p, w'u_n)).
		\]
		Then applying the claim for $n = 0$ again, we find 
		\[ 
		\delta'((\delta(p, w'), \delta(p, w'u_n)), w_n) = (\delta(p, w'u_n w_n), \delta(p, w'u_n w_n)).
		\]
		Thus we have
		$
		\applytrans{\delta'}{(q,p)}{w} = (\applytrans{\delta}{p}{w}, \applytrans{\delta}{p}{w}).
		$		
		\end{proofitem}		
	
		\item \label{constr_sync:item:QxQ}  For $v \in L \cap \Sigma^+$, 
		we have  $\delta'(Q', v) = \{(q,q) \mid q \in \delta(Q, v) \}$.
		
		\begin{proofitem}{\ref{constr_sync:item:QxQ}}
		By assumption on~$L$,
		we have $L \cap \Sigma^+ \subseteq \Sigma^*( \Sigma \setminus \Gamma )$
		and so $v \in \Sigma^* (\Sigma \setminus \Gamma)$.
		Thus, the claim follows from Property~\ref{constr_sync:item:Gamma_suffix}.		
		\end{proofitem}
	\end{enumerate}	
	
	\medskip 
	\noindent\underline{Verification:} We can assume $|Q| \ge 2$.
	First, suppose $|\delta(Q, v) | = 1$ for some $v \in L$.
	Then, by Property~\ref{constr_sync:item:QxQ}, $|\delta'(Q \times Q, v)| = 1$.
	So, $|\delta'(Q', vu)| = 1$ for $u \in U$.	
	
	Conversely, assume $|\delta'(Q', vu)| = 1$
	for $v \in L$ and $u \in U$. %, where we can assume $|v| > 0$.
%	By Property~\ref{constr_sync:item:Gamma_suffix}, as $v \notin \Sigma^*\Gamma$,
%	we have $\delta'(Q\times Q, v) = \{(q,q) \mid q \in \delta(Q, v) \}$.
	By Property~\ref{constr_sync:item:word_in_Gamma}, for $p,q \in Q$ and $u \in \Gamma^*$,
	we have $\delta'((q,q), u) = (q, \delta(q, u))$ and $\delta'((p,p), u) = (p, \delta(p, u))$.
	So,
	\[
	\delta'((q,q), u) = \delta'((p,p), u)
	\Leftrightarrow 
	q = p.
	\]
	Hence, with Property~\ref{constr_sync:item:QxQ}, $|\delta'(Q', vu)| = 1$
	implies $|\delta'(Q', v)| = 1$,
	which, again by Property~\ref{constr_sync:item:QxQ}, then implies $|\delta(Q, v)| = 1$.	
%	
%	
%	
%	
%	
%	
%	\bigskip 
%	
%	Let $v \in L$ and $u \in U$. We can assume $|v| > 0$.
%	By Property~\ref{constr_sync:item:Gamma_suffix}, as $v \notin \Sigma^*\Gamma$,
%	we have $\delta'(Q\times Q, v) = \{(q,q) \mid q \in \delta(Q, v) \}$.
%	So,  $|\delta(Q, v) | = 1$,
%	implies $|\delta'(Q \times Q, vu)| = 1$.
%	
%	%By Property~\ref{} for $v \in \Gamma^*$ 
%	%we have $\delta'(S \times T, v) = \bigcup_{q \in S} \{q\} \times \delta(T, v)$.
%	
%	Furthermore, by Property~\ref{constr_sync:item:word_in_Gamma}, as $u \in \Gamma^*$, we have $(q, \delta(q, u)) = \delta'((q,q), u)$ and $(p, \delta(p, u)) = \delta'((p,p), u)$,
%	and so $\delta'((q,q), u) = \delta'((p,p), u) = (p, \delta(u, p))$
%	if and only if $q = p$ for $q,p \in Q$. This yields
%	that $|\delta'(Q \times Q, vu)| = 1$
%	implies $|\delta'(Q \times Q, v)| = 1$,
%	which then implies $|\delta(Q, v)| = 1$.
\end{proof} 

Lastly, we give a condition for preservation of the complexity w.r.t.~concatenation.

\begin{proposition} \label{constrsync:genresults:prop:concat}
	Let $L_1, L_2, L_3 \subseteq \Sigma^*$ be regular languages.
	If $L_1 L_2 L_3 \subseteq \fact{L_2}$, then $
	\constrsync{L_1 L_2 L_3} \eqACZEROm \constrsync{L_2}.
	$
\end{proposition}
\begin{proof}
	In general, it is $\fact{L_2} 
	\subseteq 
	\fact{L_1 L_2 L_3}$.
	By assumption, $\fact{L_1 L_2 L_3}
	\subseteq \fact{L_2}$.
	So, the claim follows from~\cref{gen_results:cor:factors_etc_complexity}.
\end{proof}

\begin{corollary}
 	For $u,v \in \Sigma^*$ and $L \subseteq \Sigma^*$ with $uLv \subseteq \fact{L}$
 	we have $\constrsync{uLv} \eqACZEROm \constrsync{L}$.
\end{corollary}